% ============================================================================
%  figura1_tikz.tex  -  Figura 1 (versión 2): estructura del circuito
%  Pong en FPGA - Electrónica Digital 1 - UNAL
%  Daniel Felipe Castro González, Juan David Ortiz Muñoz, David Eduardo Chivata Arias
% ----------------------------------------------------------------------------
%  Solo trae el tikzpicture, para poder meterlo tal cual en la guía con
%  % ============================================================================
%  figura1_tikz.tex  -  Figura 1 (versión 2): estructura del circuito
%  Pong en FPGA - Electrónica Digital 1 - UNAL
%  Daniel Felipe Castro González, Juan David Ortiz Muñoz, David Eduardo Chivata Arias
% ----------------------------------------------------------------------------
%  Solo trae el tikzpicture, para poder meterlo tal cual en la guía con
%  % ============================================================================
%  figura1_tikz.tex  -  Figura 1 (versión 2): estructura del circuito
%  Pong en FPGA - Electrónica Digital 1 - UNAL
%  Daniel Felipe Castro González, Juan David Ortiz Muñoz, David Eduardo Chivata Arias
% ----------------------------------------------------------------------------
%  Solo trae el tikzpicture, para poder meterlo tal cual en la guía con
%  % ============================================================================
%  figura1_tikz.tex  -  Figura 1 (versión 2): estructura del circuito
%  Pong en FPGA - Electrónica Digital 1 - UNAL
%  Daniel Felipe Castro González, Juan David Ortiz Muñoz, David Eduardo Chivata Arias
% ----------------------------------------------------------------------------
%  Solo trae el tikzpicture, para poder meterlo tal cual en la guía con
%  \input{figura1_tikz.tex}. En el preámbulo de la guía hay que tener:
%     \usepackage{tikz}
%     \usetikzlibrary{positioning,arrows.meta,calc,fit,backgrounds}
%  Los colores marcan en qué etapa del paso a paso se arma cada bloque.
% ============================================================================
\begin{tikzpicture}[
    font=\footnotesize,
    >={Stealth[length=2.2mm]},
    blk/.style={draw, rounded corners=2pt, align=center, minimum width=3.1cm,
                minimum height=1.0cm, inner sep=3pt, line width=0.6pt},
    ext/.style={draw, align=center, minimum width=2.6cm, minimum height=0.9cm,
                inner sep=3pt, fill=gray!8, line width=0.5pt},
    e1/.style={blk, fill=blue!12},
    e2/.style={blk, fill=green!14},
    e3/.style={blk, fill=orange!18},
    e4/.style={blk, fill=violet!12},
    sig/.style={->, line width=0.6pt},
    lbl/.style={font=\scriptsize, align=center, inner sep=1.5pt},
    et/.style={font=\scriptsize\bfseries, anchor=north east, inner sep=2pt},
]

% ------------------------------------------------------------------ Entradas
\node[ext] (osc) at (-0.9, 6.6)  {Oscilador\\50 MHz (P23)};
\node[ext] (rf1) at (-0.9, 4.3)  {Receptor RF J1\\(HT12D)};
\node[ext] (rf2) at (-0.9, 2.6)  {Receptor RF J2\\(HT12D)};
\node[ext] (btn) at (-0.9, 0.9)  {K2 START (P90)\\K3 RESET (P91)};
\node[ext] (dip) at (-0.9, -0.6) {DIP SW1 (sonido)\\SW2 (calibración)};

% ------------------------------------------------------------------ FPGA
\node[e1, minimum width=2.7cm] (tck) at (4.6, 6.6) {Generador de\\bases de tiempo};
\node[e2, minimum width=2.7cm, minimum height=6.0cm] (ent) at (4.6, 2.0)
      {\textbf{Entradas}\\[6pt]sincronizador\\(2 flip-flops)\\[6pt]antirrebote\\(8 muestras\\a 1 kHz)\\[6pt]validación\\con VT};

\node[e2] (pal) at (9.4, 5.0) {Paletas J1 y J2\\(registros YP)};
\node[e3] (pel) at (9.4, 3.2) {Pelota\\+ colisiones};
\node[e4, minimum height=1.3cm, minimum width=3.5cm] (ctl) at (9.4, 1.3)
      {Control del juego (FSM)\\ESPERA, JUEGO, PUNTO,\\PAUSA, FIN};
\node[e4] (pun) at (9.4, -0.6) {Puntaje\\(2 contadores 0–9)};

\node[e1, minimum width=2.8cm] (img) at (14.2, 5.0) {Generador de\\imagen 8×32};
\node[e1, minimum width=2.8cm] (tx)  at (14.2, 3.2) {FSM y transmisor\\serial MAX7219};
\node[e4, minimum width=2.8cm] (son) at (14.2, 1.3) {Sonido\\(tonos)};
\node[e4, minimum width=2.8cm] (dmx) at (14.2, -0.6) {Display\\multiplexado};

% Etiquetas de etapa
\foreach \n/\e in {tck/E1, img/E1, tx/E1, ent/E2, pal/E2, pel/E3, ctl/E4, pun/E4, dmx/E4, son/E4}
    \node[et] at (\n.north east) {\e};

% ------------------------------------------------------------------ Salidas
\node[ext, minimum width=2.9cm, minimum height=1.3cm] (mat) at (19.9, 5.0) {Módulo MAX7219 ×4\\4 matrices 8×8\\(pantalla 8×32)};
\node[ext, minimum width=2.9cm] (buf) at (19.9, 3.2)  {SN74AHCT125\\3,3 V $\rightarrow$ 5 V};
\node[ext, minimum width=2.9cm] (buz) at (19.9, 1.3)  {Buzzer (P141)};
\node[ext, minimum width=2.9cm] (dsp) at (19.9, -0.6) {Display de 8 dígitos\\de la tarjeta};

% ------------------------------------------------------------------ Caja de la FPGA
\begin{scope}[on background layer]
  \node[draw, dashed, rounded corners=4pt, line width=0.7pt, inner sep=8pt,
        fit=(tck)(ent)(pal)(pun)(img)(dmx),
        label={[font=\small\bfseries, anchor=south west]north west:FPGA Cyclone IV E EP4CE10E22C8N (tarjeta A-C4E6E10)}] (fpga) {};
\end{scope}

% ------------------------------------------------------------------ Conexiones
% Reloj
\draw[sig] (osc) -- node[lbl, above] {CLK\_SYS} (tck);
\node[lbl, align=left, anchor=west] at ($(tck.east)+(0.2,0)$)
     {tick\_1ms, tick\_bola, parpadeo\\$\rightarrow$ a todos los bloques};

% Entradas físicas
\draw[sig] (rf1.east) -- node[lbl, above] {UP, DOWN, VT} (rf1.east -| ent.west);
\draw[sig] (rf2.east) -- node[lbl, above] {UP, DOWN, VT} (rf2.east -| ent.west);
\draw[sig] (btn.east) -- node[lbl, above] {START, RESET} (btn.east -| ent.west);
\draw[sig] (dip.east) -- node[lbl, above] {SW1, SW2} (dip.east -| ent.west);

% Entradas -> motor
\draw[sig] (ent.east |- pal.west) -- node[lbl, above] {subir, bajar\\(×2 jugadores)} (pal.west);
\draw[sig] (ent.east |- ctl.west) -- node[lbl, above] {start, reset,\\enlace} (ctl.west);

% Dentro del motor
\draw[sig] (pal) -- node[lbl, right] {YP1, YP2 [2:0]} (pel);
\draw[sig] ([xshift=-7mm]pel.south) -- node[lbl, left] {gol,\\rebote} ([xshift=-7mm]ctl.north);
\draw[sig] ([xshift=7mm]ctl.north) -- node[lbl, right] {saque,\\nivel[1:0]} ([xshift=7mm]pel.south);
\draw[sig] ([xshift=-7mm]ctl.south) -- node[lbl, left] {contar,\\limpiar} ([xshift=-7mm]pun.north);
\draw[sig] ([xshift=7mm]pun.north) -- node[lbl, right] {gana\_j1,\\gana\_j2} ([xshift=7mm]ctl.south);
\draw[sig] ([yshift=-2.5mm]pel.west) -- ++(-0.45,0) |- node[lbl, pos=0.4, left] {gol\_j1,\\gol\_j2} ([yshift=-2.5mm]pun.west);

% Motor -> visualización
\draw[sig] (pal) -- node[lbl, above] {YP1, YP2} (img);
\draw[sig] (pel.east) -- ++(0.55,0) |- ([yshift=-3mm]img.west);
\node[lbl, anchor=west] at ($(pel.east)+(0.6,0.75)$) {X[4:0],\\Y[2:0]};
\draw[sig] (img) -- node[lbl, right] {bytes[255:0]} (tx);
\draw[sig] (ctl) -- node[lbl, above] {eventos} (son);
\draw[sig] ([yshift=-4mm]ctl.east) -- ++(0.9,0) |- ([yshift=2.5mm]dmx.west);
\node[lbl, anchor=west] at ($(ctl.east)+(0.95,-1.05)$) {estado,\\nivel};
\draw[sig] ([yshift=-2.5mm]pun.east) -- node[lbl, below] {P1, P2 [3:0]} ([yshift=-2.5mm]dmx.west);

% Salidas físicas
\draw[sig] (tx) -- node[lbl, above] {DIN, CLK\_LED,\\LOAD} (buf);
\draw[sig] (buf) -- node[lbl, right] {5 V} (mat);
\draw[sig] (son) -- node[lbl, above] {BUZZER} (buz);
\draw[sig] (dmx) -- node[lbl, above] {SEG[7:0],\\DIG[7:0]} (dsp);

% ------------------------------------------------------------------ Leyenda
\matrix[anchor=north west, column sep=4pt, row sep=0pt, font=\scriptsize]
  at ($(fpga.south west)+(0,-0.25)$) {
  \node[e1, minimum width=0.5cm, minimum height=0.3cm] {}; & \node[anchor=west]{E1: pantalla}; &
  \node[e2, minimum width=0.5cm, minimum height=0.3cm] {}; & \node[anchor=west]{E2: entradas y paletas}; &
  \node[e3, minimum width=0.5cm, minimum height=0.3cm] {}; & \node[anchor=west]{E3: pelota y colisiones}; &
  \node[e4, minimum width=0.5cm, minimum height=0.3cm] {}; & \node[anchor=west]{E4: reglas, marcador y sonido}; &
  \node[ext, minimum width=0.5cm, minimum height=0.3cm] {}; & \node[anchor=west]{fuera de la FPGA}; \\
};
\end{tikzpicture}
. En el preámbulo de la guía hay que tener:
%     \usepackage{tikz}
%     \usetikzlibrary{positioning,arrows.meta,calc,fit,backgrounds}
%  Los colores marcan en qué etapa del paso a paso se arma cada bloque.
% ============================================================================
\begin{tikzpicture}[
    font=\footnotesize,
    >={Stealth[length=2.2mm]},
    blk/.style={draw, rounded corners=2pt, align=center, minimum width=3.1cm,
                minimum height=1.0cm, inner sep=3pt, line width=0.6pt},
    ext/.style={draw, align=center, minimum width=2.6cm, minimum height=0.9cm,
                inner sep=3pt, fill=gray!8, line width=0.5pt},
    e1/.style={blk, fill=blue!12},
    e2/.style={blk, fill=green!14},
    e3/.style={blk, fill=orange!18},
    e4/.style={blk, fill=violet!12},
    sig/.style={->, line width=0.6pt},
    lbl/.style={font=\scriptsize, align=center, inner sep=1.5pt},
    et/.style={font=\scriptsize\bfseries, anchor=north east, inner sep=2pt},
]

% ------------------------------------------------------------------ Entradas
\node[ext] (osc) at (-0.9, 6.6)  {Oscilador\\50 MHz (P23)};
\node[ext] (rf1) at (-0.9, 4.3)  {Receptor RF J1\\(HT12D)};
\node[ext] (rf2) at (-0.9, 2.6)  {Receptor RF J2\\(HT12D)};
\node[ext] (btn) at (-0.9, 0.9)  {K2 START (P90)\\K3 RESET (P91)};
\node[ext] (dip) at (-0.9, -0.6) {DIP SW1 (sonido)\\SW2 (calibración)};

% ------------------------------------------------------------------ FPGA
\node[e1, minimum width=2.7cm] (tck) at (4.6, 6.6) {Generador de\\bases de tiempo};
\node[e2, minimum width=2.7cm, minimum height=6.0cm] (ent) at (4.6, 2.0)
      {\textbf{Entradas}\\[6pt]sincronizador\\(2 flip-flops)\\[6pt]antirrebote\\(8 muestras\\a 1 kHz)\\[6pt]validación\\con VT};

\node[e2] (pal) at (9.4, 5.0) {Paletas J1 y J2\\(registros YP)};
\node[e3] (pel) at (9.4, 3.2) {Pelota\\+ colisiones};
\node[e4, minimum height=1.3cm, minimum width=3.5cm] (ctl) at (9.4, 1.3)
      {Control del juego (FSM)\\ESPERA, JUEGO, PUNTO,\\PAUSA, FIN};
\node[e4] (pun) at (9.4, -0.6) {Puntaje\\(2 contadores 0–9)};

\node[e1, minimum width=2.8cm] (img) at (14.2, 5.0) {Generador de\\imagen 8×32};
\node[e1, minimum width=2.8cm] (tx)  at (14.2, 3.2) {FSM y transmisor\\serial MAX7219};
\node[e4, minimum width=2.8cm] (son) at (14.2, 1.3) {Sonido\\(tonos)};
\node[e4, minimum width=2.8cm] (dmx) at (14.2, -0.6) {Display\\multiplexado};

% Etiquetas de etapa
\foreach \n/\e in {tck/E1, img/E1, tx/E1, ent/E2, pal/E2, pel/E3, ctl/E4, pun/E4, dmx/E4, son/E4}
    \node[et] at (\n.north east) {\e};

% ------------------------------------------------------------------ Salidas
\node[ext, minimum width=2.9cm, minimum height=1.3cm] (mat) at (19.9, 5.0) {Módulo MAX7219 ×4\\4 matrices 8×8\\(pantalla 8×32)};
\node[ext, minimum width=2.9cm] (buf) at (19.9, 3.2)  {SN74AHCT125\\3,3 V $\rightarrow$ 5 V};
\node[ext, minimum width=2.9cm] (buz) at (19.9, 1.3)  {Buzzer (P141)};
\node[ext, minimum width=2.9cm] (dsp) at (19.9, -0.6) {Display de 8 dígitos\\de la tarjeta};

% ------------------------------------------------------------------ Caja de la FPGA
\begin{scope}[on background layer]
  \node[draw, dashed, rounded corners=4pt, line width=0.7pt, inner sep=8pt,
        fit=(tck)(ent)(pal)(pun)(img)(dmx),
        label={[font=\small\bfseries, anchor=south west]north west:FPGA Cyclone IV E EP4CE10E22C8N (tarjeta A-C4E6E10)}] (fpga) {};
\end{scope}

% ------------------------------------------------------------------ Conexiones
% Reloj
\draw[sig] (osc) -- node[lbl, above] {CLK\_SYS} (tck);
\node[lbl, align=left, anchor=west] at ($(tck.east)+(0.2,0)$)
     {tick\_1ms, tick\_bola, parpadeo\\$\rightarrow$ a todos los bloques};

% Entradas físicas
\draw[sig] (rf1.east) -- node[lbl, above] {UP, DOWN, VT} (rf1.east -| ent.west);
\draw[sig] (rf2.east) -- node[lbl, above] {UP, DOWN, VT} (rf2.east -| ent.west);
\draw[sig] (btn.east) -- node[lbl, above] {START, RESET} (btn.east -| ent.west);
\draw[sig] (dip.east) -- node[lbl, above] {SW1, SW2} (dip.east -| ent.west);

% Entradas -> motor
\draw[sig] (ent.east |- pal.west) -- node[lbl, above] {subir, bajar\\(×2 jugadores)} (pal.west);
\draw[sig] (ent.east |- ctl.west) -- node[lbl, above] {start, reset,\\enlace} (ctl.west);

% Dentro del motor
\draw[sig] (pal) -- node[lbl, right] {YP1, YP2 [2:0]} (pel);
\draw[sig] ([xshift=-7mm]pel.south) -- node[lbl, left] {gol,\\rebote} ([xshift=-7mm]ctl.north);
\draw[sig] ([xshift=7mm]ctl.north) -- node[lbl, right] {saque,\\nivel[1:0]} ([xshift=7mm]pel.south);
\draw[sig] ([xshift=-7mm]ctl.south) -- node[lbl, left] {contar,\\limpiar} ([xshift=-7mm]pun.north);
\draw[sig] ([xshift=7mm]pun.north) -- node[lbl, right] {gana\_j1,\\gana\_j2} ([xshift=7mm]ctl.south);
\draw[sig] ([yshift=-2.5mm]pel.west) -- ++(-0.45,0) |- node[lbl, pos=0.4, left] {gol\_j1,\\gol\_j2} ([yshift=-2.5mm]pun.west);

% Motor -> visualización
\draw[sig] (pal) -- node[lbl, above] {YP1, YP2} (img);
\draw[sig] (pel.east) -- ++(0.55,0) |- ([yshift=-3mm]img.west);
\node[lbl, anchor=west] at ($(pel.east)+(0.6,0.75)$) {X[4:0],\\Y[2:0]};
\draw[sig] (img) -- node[lbl, right] {bytes[255:0]} (tx);
\draw[sig] (ctl) -- node[lbl, above] {eventos} (son);
\draw[sig] ([yshift=-4mm]ctl.east) -- ++(0.9,0) |- ([yshift=2.5mm]dmx.west);
\node[lbl, anchor=west] at ($(ctl.east)+(0.95,-1.05)$) {estado,\\nivel};
\draw[sig] ([yshift=-2.5mm]pun.east) -- node[lbl, below] {P1, P2 [3:0]} ([yshift=-2.5mm]dmx.west);

% Salidas físicas
\draw[sig] (tx) -- node[lbl, above] {DIN, CLK\_LED,\\LOAD} (buf);
\draw[sig] (buf) -- node[lbl, right] {5 V} (mat);
\draw[sig] (son) -- node[lbl, above] {BUZZER} (buz);
\draw[sig] (dmx) -- node[lbl, above] {SEG[7:0],\\DIG[7:0]} (dsp);

% ------------------------------------------------------------------ Leyenda
\matrix[anchor=north west, column sep=4pt, row sep=0pt, font=\scriptsize]
  at ($(fpga.south west)+(0,-0.25)$) {
  \node[e1, minimum width=0.5cm, minimum height=0.3cm] {}; & \node[anchor=west]{E1: pantalla}; &
  \node[e2, minimum width=0.5cm, minimum height=0.3cm] {}; & \node[anchor=west]{E2: entradas y paletas}; &
  \node[e3, minimum width=0.5cm, minimum height=0.3cm] {}; & \node[anchor=west]{E3: pelota y colisiones}; &
  \node[e4, minimum width=0.5cm, minimum height=0.3cm] {}; & \node[anchor=west]{E4: reglas, marcador y sonido}; &
  \node[ext, minimum width=0.5cm, minimum height=0.3cm] {}; & \node[anchor=west]{fuera de la FPGA}; \\
};
\end{tikzpicture}
. En el preámbulo de la guía hay que tener:
%     \usepackage{tikz}
%     \usetikzlibrary{positioning,arrows.meta,calc,fit,backgrounds}
%  Los colores marcan en qué etapa del paso a paso se arma cada bloque.
% ============================================================================
\begin{tikzpicture}[
    font=\footnotesize,
    >={Stealth[length=2.2mm]},
    blk/.style={draw, rounded corners=2pt, align=center, minimum width=3.1cm,
                minimum height=1.0cm, inner sep=3pt, line width=0.6pt},
    ext/.style={draw, align=center, minimum width=2.6cm, minimum height=0.9cm,
                inner sep=3pt, fill=gray!8, line width=0.5pt},
    e1/.style={blk, fill=blue!12},
    e2/.style={blk, fill=green!14},
    e3/.style={blk, fill=orange!18},
    e4/.style={blk, fill=violet!12},
    sig/.style={->, line width=0.6pt},
    lbl/.style={font=\scriptsize, align=center, inner sep=1.5pt},
    et/.style={font=\scriptsize\bfseries, anchor=north east, inner sep=2pt},
]

% ------------------------------------------------------------------ Entradas
\node[ext] (osc) at (-0.9, 6.6)  {Oscilador\\50 MHz (P23)};
\node[ext] (rf1) at (-0.9, 4.3)  {Receptor RF J1\\(HT12D)};
\node[ext] (rf2) at (-0.9, 2.6)  {Receptor RF J2\\(HT12D)};
\node[ext] (btn) at (-0.9, 0.9)  {K2 START (P90)\\K3 RESET (P91)};
\node[ext] (dip) at (-0.9, -0.6) {DIP SW1 (sonido)\\SW2 (calibración)};

% ------------------------------------------------------------------ FPGA
\node[e1, minimum width=2.7cm] (tck) at (4.6, 6.6) {Generador de\\bases de tiempo};
\node[e2, minimum width=2.7cm, minimum height=6.0cm] (ent) at (4.6, 2.0)
      {\textbf{Entradas}\\[6pt]sincronizador\\(2 flip-flops)\\[6pt]antirrebote\\(8 muestras\\a 1 kHz)\\[6pt]validación\\con VT};

\node[e2] (pal) at (9.4, 5.0) {Paletas J1 y J2\\(registros YP)};
\node[e3] (pel) at (9.4, 3.2) {Pelota\\+ colisiones};
\node[e4, minimum height=1.3cm, minimum width=3.5cm] (ctl) at (9.4, 1.3)
      {Control del juego (FSM)\\ESPERA, JUEGO, PUNTO,\\PAUSA, FIN};
\node[e4] (pun) at (9.4, -0.6) {Puntaje\\(2 contadores 0–9)};

\node[e1, minimum width=2.8cm] (img) at (14.2, 5.0) {Generador de\\imagen 8×32};
\node[e1, minimum width=2.8cm] (tx)  at (14.2, 3.2) {FSM y transmisor\\serial MAX7219};
\node[e4, minimum width=2.8cm] (son) at (14.2, 1.3) {Sonido\\(tonos)};
\node[e4, minimum width=2.8cm] (dmx) at (14.2, -0.6) {Display\\multiplexado};

% Etiquetas de etapa
\foreach \n/\e in {tck/E1, img/E1, tx/E1, ent/E2, pal/E2, pel/E3, ctl/E4, pun/E4, dmx/E4, son/E4}
    \node[et] at (\n.north east) {\e};

% ------------------------------------------------------------------ Salidas
\node[ext, minimum width=2.9cm, minimum height=1.3cm] (mat) at (19.9, 5.0) {Módulo MAX7219 ×4\\4 matrices 8×8\\(pantalla 8×32)};
\node[ext, minimum width=2.9cm] (buf) at (19.9, 3.2)  {SN74AHCT125\\3,3 V $\rightarrow$ 5 V};
\node[ext, minimum width=2.9cm] (buz) at (19.9, 1.3)  {Buzzer (P141)};
\node[ext, minimum width=2.9cm] (dsp) at (19.9, -0.6) {Display de 8 dígitos\\de la tarjeta};

% ------------------------------------------------------------------ Caja de la FPGA
\begin{scope}[on background layer]
  \node[draw, dashed, rounded corners=4pt, line width=0.7pt, inner sep=8pt,
        fit=(tck)(ent)(pal)(pun)(img)(dmx),
        label={[font=\small\bfseries, anchor=south west]north west:FPGA Cyclone IV E EP4CE10E22C8N (tarjeta A-C4E6E10)}] (fpga) {};
\end{scope}

% ------------------------------------------------------------------ Conexiones
% Reloj
\draw[sig] (osc) -- node[lbl, above] {CLK\_SYS} (tck);
\node[lbl, align=left, anchor=west] at ($(tck.east)+(0.2,0)$)
     {tick\_1ms, tick\_bola, parpadeo\\$\rightarrow$ a todos los bloques};

% Entradas físicas
\draw[sig] (rf1.east) -- node[lbl, above] {UP, DOWN, VT} (rf1.east -| ent.west);
\draw[sig] (rf2.east) -- node[lbl, above] {UP, DOWN, VT} (rf2.east -| ent.west);
\draw[sig] (btn.east) -- node[lbl, above] {START, RESET} (btn.east -| ent.west);
\draw[sig] (dip.east) -- node[lbl, above] {SW1, SW2} (dip.east -| ent.west);

% Entradas -> motor
\draw[sig] (ent.east |- pal.west) -- node[lbl, above] {subir, bajar\\(×2 jugadores)} (pal.west);
\draw[sig] (ent.east |- ctl.west) -- node[lbl, above] {start, reset,\\enlace} (ctl.west);

% Dentro del motor
\draw[sig] (pal) -- node[lbl, right] {YP1, YP2 [2:0]} (pel);
\draw[sig] ([xshift=-7mm]pel.south) -- node[lbl, left] {gol,\\rebote} ([xshift=-7mm]ctl.north);
\draw[sig] ([xshift=7mm]ctl.north) -- node[lbl, right] {saque,\\nivel[1:0]} ([xshift=7mm]pel.south);
\draw[sig] ([xshift=-7mm]ctl.south) -- node[lbl, left] {contar,\\limpiar} ([xshift=-7mm]pun.north);
\draw[sig] ([xshift=7mm]pun.north) -- node[lbl, right] {gana\_j1,\\gana\_j2} ([xshift=7mm]ctl.south);
\draw[sig] ([yshift=-2.5mm]pel.west) -- ++(-0.45,0) |- node[lbl, pos=0.4, left] {gol\_j1,\\gol\_j2} ([yshift=-2.5mm]pun.west);

% Motor -> visualización
\draw[sig] (pal) -- node[lbl, above] {YP1, YP2} (img);
\draw[sig] (pel.east) -- ++(0.55,0) |- ([yshift=-3mm]img.west);
\node[lbl, anchor=west] at ($(pel.east)+(0.6,0.75)$) {X[4:0],\\Y[2:0]};
\draw[sig] (img) -- node[lbl, right] {bytes[255:0]} (tx);
\draw[sig] (ctl) -- node[lbl, above] {eventos} (son);
\draw[sig] ([yshift=-4mm]ctl.east) -- ++(0.9,0) |- ([yshift=2.5mm]dmx.west);
\node[lbl, anchor=west] at ($(ctl.east)+(0.95,-1.05)$) {estado,\\nivel};
\draw[sig] ([yshift=-2.5mm]pun.east) -- node[lbl, below] {P1, P2 [3:0]} ([yshift=-2.5mm]dmx.west);

% Salidas físicas
\draw[sig] (tx) -- node[lbl, above] {DIN, CLK\_LED,\\LOAD} (buf);
\draw[sig] (buf) -- node[lbl, right] {5 V} (mat);
\draw[sig] (son) -- node[lbl, above] {BUZZER} (buz);
\draw[sig] (dmx) -- node[lbl, above] {SEG[7:0],\\DIG[7:0]} (dsp);

% ------------------------------------------------------------------ Leyenda
\matrix[anchor=north west, column sep=4pt, row sep=0pt, font=\scriptsize]
  at ($(fpga.south west)+(0,-0.25)$) {
  \node[e1, minimum width=0.5cm, minimum height=0.3cm] {}; & \node[anchor=west]{E1: pantalla}; &
  \node[e2, minimum width=0.5cm, minimum height=0.3cm] {}; & \node[anchor=west]{E2: entradas y paletas}; &
  \node[e3, minimum width=0.5cm, minimum height=0.3cm] {}; & \node[anchor=west]{E3: pelota y colisiones}; &
  \node[e4, minimum width=0.5cm, minimum height=0.3cm] {}; & \node[anchor=west]{E4: reglas, marcador y sonido}; &
  \node[ext, minimum width=0.5cm, minimum height=0.3cm] {}; & \node[anchor=west]{fuera de la FPGA}; \\
};
\end{tikzpicture}
. En el preámbulo de la guía hay que tener:
%     \usepackage{tikz}
%     \usetikzlibrary{positioning,arrows.meta,calc,fit,backgrounds}
%  Los colores marcan en qué etapa del paso a paso se arma cada bloque.
% ============================================================================
\begin{tikzpicture}[
    font=\footnotesize,
    >={Stealth[length=2.2mm]},
    blk/.style={draw, rounded corners=2pt, align=center, minimum width=3.1cm,
                minimum height=1.0cm, inner sep=3pt, line width=0.6pt},
    ext/.style={draw, align=center, minimum width=2.6cm, minimum height=0.9cm,
                inner sep=3pt, fill=gray!8, line width=0.5pt},
    e1/.style={blk, fill=blue!12},
    e2/.style={blk, fill=green!14},
    e3/.style={blk, fill=orange!18},
    e4/.style={blk, fill=violet!12},
    sig/.style={->, line width=0.6pt},
    lbl/.style={font=\scriptsize, align=center, inner sep=1.5pt},
    et/.style={font=\scriptsize\bfseries, anchor=north east, inner sep=2pt},
]

% ------------------------------------------------------------------ Entradas
\node[ext] (osc) at (-0.9, 6.6)  {Oscilador\\50 MHz (P23)};
\node[ext] (rf1) at (-0.9, 4.3)  {Receptor RF J1\\(HT12D)};
\node[ext] (rf2) at (-0.9, 2.6)  {Receptor RF J2\\(HT12D)};
\node[ext] (btn) at (-0.9, 0.9)  {K2 START (P90)\\K3 RESET (P91)};
\node[ext] (dip) at (-0.9, -0.6) {DIP SW1 (sonido)\\SW2 (calibración)};

% ------------------------------------------------------------------ FPGA
\node[e1, minimum width=2.7cm] (tck) at (4.6, 6.6) {Generador de\\bases de tiempo};
\node[e2, minimum width=2.7cm, minimum height=6.0cm] (ent) at (4.6, 2.0)
      {\textbf{Entradas}\\[6pt]sincronizador\\(2 flip-flops)\\[6pt]antirrebote\\(8 muestras\\a 1 kHz)\\[6pt]validación\\con VT};

\node[e2] (pal) at (9.4, 5.0) {Paletas J1 y J2\\(registros YP)};
\node[e3] (pel) at (9.4, 3.2) {Pelota\\+ colisiones};
\node[e4, minimum height=1.3cm, minimum width=3.5cm] (ctl) at (9.4, 1.3)
      {Control del juego (FSM)\\ESPERA, JUEGO, PUNTO,\\PAUSA, FIN};
\node[e4] (pun) at (9.4, -0.6) {Puntaje\\(2 contadores 0–9)};

\node[e1, minimum width=2.8cm] (img) at (14.2, 5.0) {Generador de\\imagen 8×32};
\node[e1, minimum width=2.8cm] (tx)  at (14.2, 3.2) {FSM y transmisor\\serial MAX7219};
\node[e4, minimum width=2.8cm] (son) at (14.2, 1.3) {Sonido\\(tonos)};
\node[e4, minimum width=2.8cm] (dmx) at (14.2, -0.6) {Display\\multiplexado};

% Etiquetas de etapa
\foreach \n/\e in {tck/E1, img/E1, tx/E1, ent/E2, pal/E2, pel/E3, ctl/E4, pun/E4, dmx/E4, son/E4}
    \node[et] at (\n.north east) {\e};

% ------------------------------------------------------------------ Salidas
\node[ext, minimum width=2.9cm, minimum height=1.3cm] (mat) at (19.9, 5.0) {Módulo MAX7219 ×4\\4 matrices 8×8\\(pantalla 8×32)};
\node[ext, minimum width=2.9cm] (buf) at (19.9, 3.2)  {SN74AHCT125\\3,3 V $\rightarrow$ 5 V};
\node[ext, minimum width=2.9cm] (buz) at (19.9, 1.3)  {Buzzer (P141)};
\node[ext, minimum width=2.9cm] (dsp) at (19.9, -0.6) {Display de 8 dígitos\\de la tarjeta};

% ------------------------------------------------------------------ Caja de la FPGA
\begin{scope}[on background layer]
  \node[draw, dashed, rounded corners=4pt, line width=0.7pt, inner sep=8pt,
        fit=(tck)(ent)(pal)(pun)(img)(dmx),
        label={[font=\small\bfseries, anchor=south west]north west:FPGA Cyclone IV E EP4CE10E22C8N (tarjeta A-C4E6E10)}] (fpga) {};
\end{scope}

% ------------------------------------------------------------------ Conexiones
% Reloj
\draw[sig] (osc) -- node[lbl, above] {CLK\_SYS} (tck);
\node[lbl, align=left, anchor=west] at ($(tck.east)+(0.2,0)$)
     {tick\_1ms, tick\_bola, parpadeo\\$\rightarrow$ a todos los bloques};

% Entradas físicas
\draw[sig] (rf1.east) -- node[lbl, above] {UP, DOWN, VT} (rf1.east -| ent.west);
\draw[sig] (rf2.east) -- node[lbl, above] {UP, DOWN, VT} (rf2.east -| ent.west);
\draw[sig] (btn.east) -- node[lbl, above] {START, RESET} (btn.east -| ent.west);
\draw[sig] (dip.east) -- node[lbl, above] {SW1, SW2} (dip.east -| ent.west);

% Entradas -> motor
\draw[sig] (ent.east |- pal.west) -- node[lbl, above] {subir, bajar\\(×2 jugadores)} (pal.west);
\draw[sig] (ent.east |- ctl.west) -- node[lbl, above] {start, reset,\\enlace} (ctl.west);

% Dentro del motor
\draw[sig] (pal) -- node[lbl, right] {YP1, YP2 [2:0]} (pel);
\draw[sig] ([xshift=-7mm]pel.south) -- node[lbl, left] {gol,\\rebote} ([xshift=-7mm]ctl.north);
\draw[sig] ([xshift=7mm]ctl.north) -- node[lbl, right] {saque,\\nivel[1:0]} ([xshift=7mm]pel.south);
\draw[sig] ([xshift=-7mm]ctl.south) -- node[lbl, left] {contar,\\limpiar} ([xshift=-7mm]pun.north);
\draw[sig] ([xshift=7mm]pun.north) -- node[lbl, right] {gana\_j1,\\gana\_j2} ([xshift=7mm]ctl.south);
\draw[sig] ([yshift=-2.5mm]pel.west) -- ++(-0.45,0) |- node[lbl, pos=0.4, left] {gol\_j1,\\gol\_j2} ([yshift=-2.5mm]pun.west);

% Motor -> visualización
\draw[sig] (pal) -- node[lbl, above] {YP1, YP2} (img);
\draw[sig] (pel.east) -- ++(0.55,0) |- ([yshift=-3mm]img.west);
\node[lbl, anchor=west] at ($(pel.east)+(0.6,0.75)$) {X[4:0],\\Y[2:0]};
\draw[sig] (img) -- node[lbl, right] {bytes[255:0]} (tx);
\draw[sig] (ctl) -- node[lbl, above] {eventos} (son);
\draw[sig] ([yshift=-4mm]ctl.east) -- ++(0.9,0) |- ([yshift=2.5mm]dmx.west);
\node[lbl, anchor=west] at ($(ctl.east)+(0.95,-1.05)$) {estado,\\nivel};
\draw[sig] ([yshift=-2.5mm]pun.east) -- node[lbl, below] {P1, P2 [3:0]} ([yshift=-2.5mm]dmx.west);

% Salidas físicas
\draw[sig] (tx) -- node[lbl, above] {DIN, CLK\_LED,\\LOAD} (buf);
\draw[sig] (buf) -- node[lbl, right] {5 V} (mat);
\draw[sig] (son) -- node[lbl, above] {BUZZER} (buz);
\draw[sig] (dmx) -- node[lbl, above] {SEG[7:0],\\DIG[7:0]} (dsp);

% ------------------------------------------------------------------ Leyenda
\matrix[anchor=north west, column sep=4pt, row sep=0pt, font=\scriptsize]
  at ($(fpga.south west)+(0,-0.25)$) {
  \node[e1, minimum width=0.5cm, minimum height=0.3cm] {}; & \node[anchor=west]{E1: pantalla}; &
  \node[e2, minimum width=0.5cm, minimum height=0.3cm] {}; & \node[anchor=west]{E2: entradas y paletas}; &
  \node[e3, minimum width=0.5cm, minimum height=0.3cm] {}; & \node[anchor=west]{E3: pelota y colisiones}; &
  \node[e4, minimum width=0.5cm, minimum height=0.3cm] {}; & \node[anchor=west]{E4: reglas, marcador y sonido}; &
  \node[ext, minimum width=0.5cm, minimum height=0.3cm] {}; & \node[anchor=west]{fuera de la FPGA}; \\
};
\end{tikzpicture}
